% GAL1897_W06 — English translation of physical PDF pages 042–043.
% Primary authority: the frozen 1897 French diplomatic transcription from
% JP2 leaves L0041–L0042, navigated against the 96-page 1897 PDF.
% Definite source-level mathematical errors are translated without silent repair
% and identified in adjacent critical notes from the frozen critical baseline.

\providecommand{\GalSourcePageStart}[4]{}%
\providecommand{\GalSourcePageBoundary}[4]{\space}%
\providecommand{\GalSourceError}[2]{#1}%
\providecommand{\GalHistoricalFootnote}[1]{\footnote{#1}}%
\providecommand{\GalArticleEndRule}{\par\smallskip\begin{center}\rule{0.18\linewidth}{0.4pt}\end{center}}%

\GalSourcePageStart{042}{L0041}{13}{13}

% ENSEG:W06:PDF042:0001
\begin{center}
{\Large\bfseries NOTE\par}
\vspace{0.15em}
{\scriptsize\bfseries ON THE\par}
\vspace{0.18em}
{\normalsize\bfseries SOLUTION OF NUMERICAL EQUATIONS\,\textsuperscript{(*)}.\par}
\vspace{0.7em}
\rule{0.12\linewidth}{0.4pt}
\end{center}

% ENSEG:W06:PDF042:0002
\GalHistoricalFootnote{M.~Férussac's \emph{Bulletin des Sciences mathématiques}, vol.~XIII,
p.~413 (1830, June issue). \textsc{(J. Liouville.)}}

% ENSEG:W06:PDF042:0003
M.~Legendre was the first to observe that, when an algebraic equation
was put in the form
% ENSEG:W06:PDF042:0004
\[
\varphi x=x,
\]
% ENSEG:W06:PDF042:0005
where $\varphi x$ is a function of $x$ that increases continually with
$x$, it was easy to find the root of this equation immediately below
a given number $a$ if $\varphi a<a$, and the root immediately above $a$
if $\varphi a>a$.

% ENSEG:W06:PDF042:0006
To prove this, construct the curve $y=\varphi x$ and the line $y=x$.
Take an abscissa equal to $a$, and suppose, to fix ideas, that
$\varphi a>a$. I say that it will be easy to obtain the root immediately
above $a$. Indeed, the roots of the equation $\varphi x=x$ are simply
the abscissas of the points where the line and the curve intersect, and
it is clear that one will approach the nearest point of intersection by
replacing the abscissa $a$ with the abscissa $\varphi a$. A still closer
value will be obtained by taking $\varphi\varphi a$, then
$\varphi\varphi\varphi a$, and so on.

% ENSEG:W06:PDF042:0007
Let $Fx=0$ be a given equation of degree $n$, and let $Fx=X-Y$,
where $X$ and $Y$ contain only positive terms. Legendre successively
puts the equation into these two forms:
% ENSEG:W06:PDF042:0008
\[
x=\varphi x=\sqrt[n]{\frac{X}{\left(\dfrac{Y}{x^n}\right)}},
\qquad
\GalSourceError{x=\psi x=\sqrt[n]{\frac{X}{\left(\dfrac{x^n}{Y}\right)}}}{W06-PE001: the printed second radical is not in general equivalent to F(x)=0 and need not be the opposite-side fixed-point map; preserve it and consult the separate exact certificate};
\]
% ENSEG:W06:PDF042:0009
the two functions $\varphi x$ and $\psi x$ are always, as one sees,
one greater and the other less than $x$. Thus, with the aid of these two
functions, one can obtain the two roots of the equation nearest to a
given number $a$, one above and the other below it.

% ENSEG:W06:PDF042:0010
\GalCriticalNote{W06-PE001}{The second iteration formula is not equivalent to $F(x)=0$}{The
printed map, with numerator $X$ and denominator $x^n/Y$, is not generally
equivalent to $F(x)=0$ and need not lie on the side of $x$ opposite the first
map. The symmetric candidate
$x=\psi x=\sqrt[n]{\frac{Y}{\left(\dfrac{X}{x^n}\right)}}$ is an editorial
reconstruction, not a reading found in the 1897 edition. The 1846 publication
repeats the printed form, so it establishes transmission rather than a
correction. Numerical use therefore requires the reconstruction to be stated
explicitly.}

\GalSourcePageBoundary{043}{L0042}{14}{14}

% ENSEG:W06:PDF043:0011
But this method has the disadvantage of requiring the extraction of an
$n$th root at each operation. Here are two more convenient forms. Let us
seek a number $k$ such that the function
% ENSEG:W06:PDF043:0012
\[
x+\frac{Fx}{kx^n}
\]
% ENSEG:W06:PDF043:0013
increases with $x$ when $x>1$. (It is enough, in fact, to know how to
find the roots of an equation that are greater than unity.)

% ENSEG:W06:PDF043:0014
For the proposed condition we shall have
% ENSEG:W06:PDF043:0015
\[
1+\frac{d\,\dfrac{X-Y}{kx^n}}{dx}>0,
\qquad\text{or}\qquad
1-\frac{nX-xX'}{kx^{n+1}}+\frac{nY-xY'}{kx^{n+1}}>0;
\]
% ENSEG:W06:PDF043:0016
but identically we have
% ENSEG:W06:PDF043:0017
\[
\GalSourceError{nX-xX'>0,\qquad nY-xY'>0}{W06-SE001: strict positivity is false for admissible positive-term decompositions, and the prescribed value k=(nX-xX')|_{x=1} can therefore be zero};
\]
% ENSEG:W06:PDF043:0018
it is therefore enough to set
% ENSEG:W06:PDF043:0019
\[
\frac{nX-xX'}{kx^{n+1}}<1\qquad\text{for}\qquad x>1,
\]
% ENSEG:W06:PDF043:0020
and for this it is enough to take as $k$ the value of the function
$\GalSourceError{nX-xX'}{W06-SE001}$ corresponding to $x=1$.

% ENSEG:W06:PDF043:0021
\GalCriticalNote{W06-SE001}{Strict positivity requires a lower-degree term}{For
a polynomial with positive terms, $X=\sum c_jx^j$, one has
$nX-xX'=\sum_{j<n}(n-j)c_jx^j\geq0$. It is strictly positive only if some
coefficient below degree $n$ is positive; when $X=x^n$, it vanishes identically
and the prescribed $k$ is zero. The same qualification is required for $Y$.
The convergence claim therefore needs these nonvanishing hypotheses.}

% ENSEG:W06:PDF043:0022
One similarly finds a number $h$ such that the function
% ENSEG:W06:PDF043:0023
\[
x-\frac{Fx}{hx^n}
\]
% ENSEG:W06:PDF043:0024
will increase with $x$ when $x>1$, by replacing $Y$ with $X$.

% ENSEG:W06:PDF043:0025
Thus the given equation can be put into either of the forms
% ENSEG:W06:PDF043:0026
\[
x=x+\frac{Fx}{kx^n},
\qquad
x=x-\frac{Fx}{hx^n},
\]
% ENSEG:W06:PDF043:0027
both of which are rational and provide an easy method of solution.

\GalArticleEndRule
